\باب{سمسن کا ایک  تہائی کا   قاعدہ}\شناخت{ضمیمہ_د}

تکمل \عددی{\int_a^b f(x)\dif x} کی تخمین کا سمسن قاعدہ \عددی{f} کی ترسیم کو مکافی      قوسوں سے تخمین دینے پر منحصر ہے۔

\begin{figure}
\begin{minipage}{0.45\textwidth}
\centering
\begin{tikzpicture}[declare function={fx(\x)=\x;fy(\x)=3-(\x-0.25)^2;}]
\begin{axis}[hide y axis=true, clip=false,small,view/h=120,axis lines=middle,xlabel={$x$},ylabel={$y$},zlabel={$z$},xtick={-0.75,0.25,1.25},ytick={\empty},ztick={\empty},xticklabels={$x_0$,$x_1$,$x_2$}, xlabel style={anchor=west},ylabel style={anchor=east}]
\addplot[thick,domain=-1:1.5]({fx(\x)},{fy(\x)})node[pos=0.35,pin=135:{\text{\RL{قطع مکافی}}}]{};
\addplot[] coordinates {({fx(-0.75)},{fy(-0.75)}) ({fx(-0.75)},{0})}node[pos=0,left,yshift=-0.5em]{$(x_0,y_0)$};
\addplot[] coordinates {({fx(0.25)},{fy(0.25)}) ({fx(0.25)},{0})}node[pos=0,above]{$(x_1,y_1)$};
\addplot[] coordinates {({fx(1.25)},{fy(1.25)}) ({fx(1.25)},{0})}node[pos=0,right]{$(x_2,y_2)$};
\addplot[name path=kcur, thick,smooth,tension=0.4]coordinates{(-1.5,2.5)({fx(-0.75)},{fy(-0.75)})  ({fx(0.25)},{fy(0.25)}) (0.75,3) ({fx(1.25)},{fy(1.25)}) (2,-0.5) (5,-1.25)};
\addplot[]coordinates{(-0.25,0)}node[above]{$h$};
\addplot[]coordinates{(0.75,0)}node[above]{$h$};
\path[name path=kA] (2.25,0)--(2.25,-1.25);
\path[name path=kB] (3.25,0)--(3.25,-1.25);
\path[name path=kC] (4.25,0)--(4.25,-1.25);
\draw[name intersections={of={kcur and kA}}] (2.25,0)--(intersection-1)node[below,yshift=-0.5em]{$y=f(x)$};
\draw[name intersections={of={kcur and kB}}] (3.25,0)--(intersection-1);
\draw[name intersections={of={kcur and kC}}] (4.25,0)--(intersection-1);
\addplot[name path=kcur,draw=none,mark=none,domain=-0.75:1.25]({fx(\x)},{fy(\x)});
\addplot[draw=none,mark=none,name path=kaxis]coordinates {(-0.75,0)(1.25,0)};
\addplot[llgray] fill between [of=kcur and kaxis];
\end{axis}
\end{tikzpicture}
\caption{سمسن کا قاعدہ منحنی کے چھوٹے  چھوٹے قطعات کو قطع مکافی سے تخمین دیتا ہے۔}
\label{شکل_ضمیمہ_چار_اے_نو}
\end{minipage}\hfill
\begin{minipage}{0.45\textwidth}
\centering
\begin{tikzpicture}[declare function={fx(\x)=\x;fy(\x)=3-(\x-0.25)^2;}]
\begin{axis}[hide y axis=true, clip=false,small,view/h=120,axis lines=middle,xlabel={$x$},ylabel={$y$},zlabel={$z$},xtick={-0.75,0.25,1.25},ytick={\empty},ztick={\empty},xticklabels={$-h$,$0$,$h$}, xlabel style={anchor=west},ylabel style={anchor=east}]
\addplot[thick,domain=-1:1.5]({fx(\x)},{fy(\x)})node[pos=0.35,pin={[align=center,pin distance={1cm}]90:{\text{\RL{قطع مکافی}}\\ $y=Ax^2+Bx+C$}}]{};
\addplot[] coordinates {({fx(-0.75)},{fy(-0.75)}) ({fx(-0.75)},{0})}node[pos=0,left]{$(-h,y_0)$}node[pos=0.5,left]{$y_0$};
\addplot[] coordinates {({fx(0.25)},{fy(0.25)}) ({fx(0.25)},{0})}node[pos=0,above]{$(0,y_1)$};
\addplot[] coordinates {({fx(1.25)},{fy(1.25)}) ({fx(1.25)},{0})}node[pos=0,right]{$(h,y_2)$}node[pos=0.5,left]{$y_2$};
\addplot[]coordinates{({0.25},{0.5*fy(-0.75)})}node[left]{$y_1$};
\addplot[name path=kcur,draw=none,mark=none,domain=-0.75:1.25]({fx(\x)},{fy(\x)});
\addplot[draw=none,mark=none,name path=kaxis]coordinates {(-0.75,0)(1.25,0)};
\addplot[llgray] fill between [of=kcur and kaxis];
\end{axis}
\end{tikzpicture}
\caption
{سایہ دار رقبہ تلاش کرنے کے لئے \عددی{-h} تا \عددی{h} تکمل لے کر \عددی{\tfrac{h}{3}(y_0+4y_1+y_2)} حاصل ہو گا۔}
\label{شکل_ضمیمہ_چار_اے_دس}
\end{minipage}
\end{figure}

شکل \حوالہ{شکل_ضمیمہ_چار_اے_نو}  میں مکافی کے نیچے سایہ دار خطے کا رقبہ درج ذیل ہے۔
\[\text{رقبہ}=\frac{h}{3}(y_0+4y_1+y_2)\]
اس کلیہ کو \اصطلاح{سمسن کا ایک تہائی کا قاعدہ }\حاشیہب{Simson's one-third rule}\فرہنگ{سمسن! ایک تہائی  کا قاعدہ}\فرہنگ{Simson's one-third rule}کہتے ہیں۔

آئیں  یہ  کلیہ  اخذ کریں۔ الجبرا سادہ رکھنے کی غرض سے شکل \حوالہ{شکل_ضمیمہ_چار_اے_دس}  میں پیش  محددی نظام  استعمال کیا جائے گا۔مکافی کے نیچے رقبہ محور \عددی{y} کے مقام پر منحصر نہیں تاہم  اس محور   کا عمودی پیمانہ برقرار رکھنا ضروری ہے۔ مکافی کی مساوات \عددی{y=Ax^2+Bx+C} روپ رکھتی ہے لہٰذا \عددی{x=-h} تا \عددی{x=h} اس کے نیچے رقبہ درج ذیل ہو گا۔
\begin{align*}
\text{رقبہ}=\int_{-h}^{h} (Ax^2+Bx+C)\dif x&=\left[\frac{Ax^3}{3}+\frac{Bx^2}{2}+Cx\right]_{-h}^{h}\\
&=\frac{2Ah^3}{3}+2Ch\\
&=\frac{h}{3}(2Ah^2+6C)
\end{align*}
منحنی نقاط \عددی{(-h,y_0)}، \عددی{0,y_1}، اور \عددی{(h,y_2)} سے بھی گزرتی ہے لہٰذا درج ذیل بھی ہو گا۔
\[y_0=Ah^2-Bh+C,\quad y_1=C,\quad y_2=Ah^2+Bh+C\]
ان مساوات سے ذیل حاصل ہوتا ہے۔
\begin{align*}
C&=y_1\\
Ah^2-Bh&=y_0-y_1\\
Ah^2+Bh&=y_2-y_1\\
2Ah^2&=y_0+y_2-2y_1
\end{align*}
رقبے کی مساوات میں \عددی{C} اور \عددی{2Ah^2} کی قیمتیں ڈال کر ذیل  حاصل ہو گا۔
\[\text{رقبہ}=\frac{h}{3}(2Ah^2+6C)=\frac{h}{3}((y_0+y_2-2y_1)+6y_1)=\frac{h}{3}(y_0+4y_1+y_2)\]
